%------------------------------------------------------------------------------%
\begin{question}
Considere a matriz aumentada
\[
  \begin{amatrix}{3}
    1  & 2  & -1 & 4  \\
    2  & 5  & 1  & 10 \\
    -1 & -1 & 2  & -1
  \end{amatrix}
\]
\vspace{-0.5\baselineskip}
\begin{enumerate}
  \item Escalone a matriz
  \item Determine o tipo de sistema
  \item Encontre a solução do sistema se ele for possível
\end{enumerate}
\end{question}

%------------------------------------------------------------------------------%
\begin{resolution}{Escalonamento}
\begin{columns}
\column{0.35\linewidth}
  \[
    \begin{amatrix}{3}
      1  & 2  & -1 & 4  \\
      2  & 5  & 1  & 10 \\
      -1 & -1 & 2  & -1
  \end{amatrix}
  \]
  \pause
\column{0.65\linewidth}
  Eliminamos os elementos abaixo do primeiro pivô
  \\ \pause
  \(
    L_2 \leftarrow L_2 - 2 L_1
  \)
  \\ \pause
  \(
    L_3 \leftarrow L_3 + L_1
  \)
\end{columns}
\vspace{-\baselineskip}
\[
  \pause
  L'_2
  \;=\;
  L_2 - 2 L_1
  \pause
  \;=\;
  \begin{amatrix}{3}
    2  & 5  & 1  & 10
  \end{amatrix}
  - 2
  \begin{amatrix}{3}
    1  & 2  & -1 & 4
  \end{amatrix}
  \pause
  \;=\;
  \begin{amatrix}{3}
    0 & 1 & 3  & 2
  \end{amatrix}
\]
\[
  \pause
  L'_3
  \;=\;
  L_3 + L_1
  \pause
  \;=\;
  \begin{amatrix}{3}
    -1 & -1 & 2  & -1
  \end{amatrix}
  +
  \begin{amatrix}{3}
    1  & 2  & -1 & 4
  \end{amatrix}
  \pause
  \;=\;
  \begin{amatrix}{3}
    0 & 1 & 1  & 3
  \end{amatrix}
\]
\pause
\[
  \begin{amatrix}{3}
    1 & 2 & -1 & 4 \\
    0 & 1 & 3  & 2 \\
    0 & 1 & 1  & 3
  \end{amatrix}
\]
\end{resolution}

%------------------------------------------------------------------------------%
\begin{resolution}{Escalonamento}
\begin{columns}
\column{0.3\linewidth}
  \[
  \begin{amatrix}{3}
    1 & 2 & -1 & 4 \\
    0 & 1 & 3  & 2 \\
    0 & 1 & 1  & 3
  \end{amatrix}
  \]
  \pause
\column{0.7\linewidth}
  Eliminamos os elementos abaixo do segundo pivô
  \\ \pause
  \(
    L_3 \leftarrow L_3 - L_2
  \)
\end{columns}
\vspace{-\baselineskip}
\[
  \pause
  L'_3
  \;=\;
  L_3 - L_2
  \pause
  \;=\;
  \begin{amatrix}{3}
    0 & 1 & 1  & 3
  \end{amatrix}
  -
  \begin{amatrix}{3}
    0 & 1 & 3  & 2
  \end{amatrix}
  \pause
  \;=\;
  \begin{amatrix}{3}
    0 & 0 & -2 & 1
  \end{amatrix}
\]
\pause
\[
  \begin{amatrix}{3}
    1 & 2 & -1 & 4 \\
    0 & 1 & 3  & 2 \\
    0 & 0 & -2 & 1
  \end{amatrix}
\]
\end{resolution}

%------------------------------------------------------------------------------%
\begin{resolution}{Classificação}
  O sistema é possível e determinado
\end{resolution}

%------------------------------------------------------------------------------%
\begin{resolution}{Substituição Retroativa}
\vspace{-\baselineskip}
\begin{columns}[t]
\column{0.24\linewidth}
\begin{align*}
  \onslide<+->{-2z & = 1            \\}
  \onslide<+->{z   & = \frac{-1}{2}}
\end{align*}
\column{0.26\linewidth}
\begin{align*}
  \onslide<+->{y + 3z            & = 2           \\}
  \onslide<+->{y + 3\frac{-1}{2} & = 2           \\}
  \onslide<+->{y - \frac{3}{2}   & = 2           \\}
  \onslide<+->{y                 & = 2 + \frac{3}{2}                   \\}
  \onslide<+->{y                 & = \frac{2\times 2}{2} + \frac{3}{2} \\}
  \onslide<+->{y                 & = \frac{7}{2}}
\end{align*}
\column{0.44\linewidth}
\begin{align*}
  \onslide<+->{x + 2y - z                      & = 4                \\}
  \onslide<+->{x + 2\frac{7}{2} - \frac{-1}{2} & = 4                \\}
  \onslide<+->{x + 7 + \frac{1}{2}             & = 4                \\}
  \onslide<+->{x                               & = -3 - \frac{1}{2} \\}
  \onslide<+->{x                               & = -\frac{3\times 2}{2} - \frac{1}{2} \\}
  \onslide<+->{x                               & = \frac{-7}{2}}
\end{align*}
\end{columns}
\end{resolution}

%------------------------------------------------------------------------------%
\begin{resolution}{Solução}
A solução do sistema é
\(
  {\bm x}
  \;=\;
  \begin{bmatrix}
    x \\ y \\ z
  \end{bmatrix}
  \;=\;
  \left[\begin{array}{r}
      -\sfrac{7}{2} \\[3mm]
       \sfrac{7}{2} \\[3mm]
      -\sfrac{1}{2}
  \end{array}\right]
\)
\end{resolution}

%------------------------------------------------------------------------------%
\endinput
